% !TeX TS-program = xelatex
% 简历 tex 骨架 —— agent 据此手写，勿自由创造结构
% 编译：node <skill>/scripts/resume.mjs build path/to/resume.tex
% 规范：contracts/latex-rules.md
%
% profile.yaml 字段 → 章节映射（只列有对应章节的字段）：
%   meta            → 标题区（姓名/联系方式/目标岗位/照片）
%   summary         → 个人总结（可选，按 JD 方向选用）
%   education       → 教育经历
%   experience      → 实习经历
%   projects        → 项目经历
%   skills          → 技能
%   awards          → 获奖（可选）
%   open_source     → 开源（可选）
%   publications    → 论文/专利（可选；措辞必须遵循 status）
%   extras          → 兜底素材（可选，视相关性与篇幅决定去留）
% 没有对应章节的字段不代表可以忽略：无章节承载时不得写进简历，
% 但仍需留在 profile.yaml 中作为已积累素材。
\documentclass{resume}
\ResumeName{<姓名>}

\usepackage{graphicx}
% Agent 必须将 <CJK_FONT> 替换为 `node scripts/resume.mjs check-env` 输出的推荐字体。
% 标题族用合成粗体，避免与正文字重相同导致层级消失。
\setCJKmainfont{<CJK_FONT>}
\setCJKsansfont{<CJK_FONT>}
\setCJKfamilyfont{hei}[AutoFakeBold=2.5]{<CJK_FONT>}
\newcommand{\heiti}{\CJKfamily{hei}\bfseries}

% 标题区（含照片）。照片路径相对本 tex，指向 profiles/<user>/photo.jpg
\newcommand{\PhotoResumeTitle}{
  \noindent
  \begin{minipage}[c]{0.80\textwidth}
    \begin{center}
      {\zihao{-2}\heiti <姓名>}\\[-0.2em]
      <电话>~|~
      \ResumeUrl{mailto:<邮箱>}{<邮箱>}~|~
      <城市>~|~<年龄> 岁~|~<性别}\\
      <POSITION_LINE>~|~<到岗时间>~|~
      \ResumeUrl{<github url>}{github.com/...}
    \end{center}
  \end{minipage}
  \hfill
  \begin{minipage}[c]{0.14\textwidth}
    \centering
    \includegraphics[width=2.0cm]{../../profiles/<user>/photo.jpg}
  \end{minipage}
  \vspace{0.2em}
}

\begin{document}

\PhotoResumeTitle

% ========== 教育经历 ==========
\section{教育经历}
\ResumeItem[学校名|学历]{学校名}[\textnormal{专业，学院\textbar} 学历][起止]
主修课程：...
\GrayText{奖项：...}

% ========== 实习/工作经历 ==========
\section{实习经历}
\ResumeItem{公司名}[\textnormal{岗位\textbar} 其他信息][起止]
\begin{itemize}
  \item 经历要点 1（按 JD 改写，突出现相关键词，不编造）
  \item 经历要点 2
\end{itemize}

% ========== 项目经历 ==========
\section{项目经历}
\ResumeItem{项目名}[角色/技术][起止]
\begin{itemize}
  \item 项目要点（最相关排前）
  \item ...
\end{itemize}

% ========== 技能 ==========
\section{技能}
\ResumeItem{分类}{技能项列表}

% ========== 个人总结（可选，取自 profile.yaml: summary） ==========
% 按 JD 方向选一条最贴合的，不要多条堆叠；一页紧张时整节省略。
% \section{个人总结}
% \GrayText{summary[].text}

% ========== 开源（可选，取自 profile.yaml: open_source） ==========
% 与 JD 强相关时保留，否则省略。
% \section{开源}
% \ResumeItem{项目名}[作者/贡献者][起止]
% desc

% ========== 论文/专利（可选，取自 profile.yaml: publications） ==========
% status 为「已发表/已接收」才可写发表；「在审/已投稿」只能写「在审」或省略。
% \section{论文}
% \GrayText{标题（本人第 X 作者），venue，year，status}

% ========== 获奖（可选） ==========
\section{获奖}
\GrayText{...}

\end{document}
